% =============================================================
% DATOS GENERALES DEL TRABAJO DE GRADO
% Edite este archivo para personalizar la portada, el encabezado
% y la modalidad del proyecto.
% =============================================================

% -------------------------------------------------------------
% Información general
% -------------------------------------------------------------

\newcommand{\ThesisTitle}{Transformación de la dirección de rechazo en modelos de lenguaje agénticos}

% Subtítulo (deje el argumento vacío si no se usa)
\newcommand{\ThesisSubtitle}{Un análisis desde la interpretabilidad mecánica. Avances de investigación}

\newcommand{\ThesisShortTitle}{Dirección de rechazo en modelos agénticos}
\newcommand{\ThesisAuthors}{Juan Esteban Espitia Leiva}

\newcommand{\ThesisProgram}{Pregrado en Ciencia de Datos}
\newcommand{\ThesisDepartment}{Departamento de Matemáticas}
\newcommand{\ThesisUniversity}{Universidad Externado de Colombia}

\newcommand{\ThesisAdvisors}{Julián Fernando Sánchez}
\newcommand{\ThesisCity}{Bogotá, Colombia}
\newcommand{\ThesisYear}{2026}

% Archivo del logo institucional dentro de images/.
\newcommand{\UniversityLogo}{logo-externado-matematicas.png}


% =============================================================
% MODALIDAD DEL PROYECTO
% =============================================================

% Seleccione UNA de las siguientes modalidades:
%
% investigacion   = Investigación
% innovacion      = Innovación
% emprendimiento = Emprendimiento
%
% Solo debe modificar el valor de la siguiente línea.

\newcommand{\ModalidadProyecto}{investigacion}


% -------------------------------------------------------------
% Nombre de la modalidad para portada y otros elementos
% -------------------------------------------------------------

\ifdefstring{\ModalidadProyecto}{investigacion}{
  \newcommand{\TipoProyecto}{Investigación}
}{
  \ifdefstring{\ModalidadProyecto}{innovacion}{
    \newcommand{\TipoProyecto}{Innovación}
  }{
    \ifdefstring{\ModalidadProyecto}{emprendimiento}{
      \newcommand{\TipoProyecto}{Emprendimiento}
    }{
      \newcommand{\TipoProyecto}{Modalidad no válida}
    }
  }
}


% -------------------------------------------------------------
% Compatibilidad con versiones anteriores de la plantilla
% -------------------------------------------------------------

% Mantiene funcionando archivos que todavía utilicen
% el comando \ProjectType.

\newcommand{\ProjectType}{\TipoProyecto}
